% ============================================================================
%  CAPÍTULO III: ALCANCE Y LÍMITES
% ============================================================================

\section{Alcance y Límites}

\subsection{Alcance}

El presente informe documenta la configuración de los servicios de red del
proyecto integrador. Las actividades realizadas incluyen:

\subsubsection{Servidor de correo electrónico}
\begin{itemize}
  \item Instalación de Postfix y Dovecot en Debian con el administrador de
        paquetes APT.
  \item Configuración de Postfix como servidor SMTP con habilitación del
        puerto 587 (submission) para clientes autenticados.
  \item Configuración de Dovecot como servidor IMAP/POP3 con autenticación
        de usuarios vía SASL.
  \item Generación de certificados SSL/TLS auto-firmados para la protección
        de las conexiones (STARTTLS).
  \item Creación de usuarios de correo (joaquin, nicolas, david) con
        estructura de buzón en formato Maildir.
  \item Configuración del servidor DNS (BIND9) con registros MX y A para
        el dominio \texttt{sudoers.lan}.
  \item Pruebas de conectividad y envío/recepción de correo entre usuarios
        del dominio.
\end{itemize}

\subsubsection{Servidor de archivos}
\begin{itemize}
  \item Instalación de Samba en Debian.
  \item Configuración de carpetas compartidas con permisos por usuario.
  \item Pruebas de acceso desde clientes Windows y Linux.
\end{itemize}

\subsection{Límites}

Los siguientes aspectos no forman parte de este informe:

\begin{itemize}
  \item Seguridad avanzada del servidor de correo (antispam, antivirus,
        listas negras).
  \item Integración con dominios de correo externos o configuración de relay
        con proveedores externos.
  \item Configuración de registros SPF, DKIM o DMARC para autenticación de
        dominio.
  \item Desarrollo detallado del servidor de archivos (corresponde a otro
        integrante del equipo).
\end{itemize}

\subsection{Entorno de trabajo}

El entorno utilizado para la configuración consta de los siguientes
componentes:

\begin{itemize}
  \item \textbf{Servidor:} máquina con Debian 12 (Bookworm) ejecutando los
        servicios de correo y archivos directamente en el sistema operativo.
  \item \textbf{Red interna:} red administrativa \texttt{192.168.0.0/24}
        (WLO1) y red de usuarios \texttt{192.168.20.0/24} (ENO1).
  \item \textbf{Servicios instalados:} BIND9 (DNS), Postfix + Dovecot
        (correo), Samba (archivos).
  \item \textbf{Resolución de nombres:} servidor DNS propio que resuelve el
        dominio \texttt{sudoers.lan} para todos los servicios.
\end{itemize}
